\scanpage{311}

\mspage{50}\mspage{50}
% archive sheet number: 310

Pour \inserted{$\operatorname{Nerf}_1(\underline{\mathrm{Ty}})=\operatorname{Fl}(\underline{\mathrm{Ty}})$, i.e. pour} les flèches, les résultats \deleted{\unsure{du}} \unsure{]} du par. précédent (détermination de $\operatorname{Fl}(\underline{\mathrm{Ty}})$) s'interprètent aussi en disant que $\operatorname{Fl}(\underline{\mathrm{Ty}})$ est le monoïde libre engendré par un système de générateurs infini
\[
f_{\varepsilon,n},\qquad(\varepsilon,n)\in\{+1,-1\}\times\mathbb N
\]
indexé par \unsure{l'ens produit|leur produit|un produit} \deleted{\unsure{que}} $\{+1,-1\}\times\unsure{\mathcal N|\mathbb N}$. Ici $f_{\varepsilon,n}$ désigne la \inserted{unique} flèche de $\mathrm{Ty}$
\[
f_{\varepsilon,\ell}:\Delta_\varepsilon\longrightarrow\Delta_\varepsilon^{\deleted{\unsure{n}}\ell}
\]
(où bien sûr l'exposant $n$ désigne une exponentiation pour la loi de composition de $\underline{\mathrm{Ty}}$
\[
\Delta_\varepsilon^\ell=\underbrace{\Delta_\varepsilon\circ\Delta_\varepsilon\circ\cdots\circ\Delta_\varepsilon}_{\ell\ \text{copies}},
\]
et non la puissance cartésienne dans $\mathrm{Cat}$, ou dans $\underline{\mathrm{Ty}}$ (où elle n'existe que si $n\in\{0,1\}$)).

\deleted{\unsure{on}} Dans la description de $\operatorname{Fl}(\underline{\mathrm{Ty}})$ dans le par. précédent, on a écrit le ``terme général'' du monoïde libre en groupant ensemble les \deleted{\unsure{f}}\inserted{groupes de} facteurs successifs de la forme $f_{+,0}, f_{-,0}$ i.e. correspondant \deleted{\unsure{aux deux}}\inserted{aux} flèches \inserted{de} dégénérescence indécomposables (pour $*$)
\[
\left\{\begin{aligned}
 f_{+,0}:{}&\Delta_+\unsure{\text{ --- }|\longrightarrow}e\\
 f_{-,0}:{}&\Delta_-\unsure{\text{ --- }|\longrightarrow}e
\end{aligned}\right.
\]
% slanted
\margin{\uline{NB} Le ``mot'' vide de $\operatorname{Fl}(\underline{\mathrm{Ty}})$ \deleted{\unsure{est}} correspond à la flèche $\mathrm{id}_1$, où $1$ est l'élément unité de $\mathrm{Ty}$, \deleted{\unsure{cor}} \unsure{représenté} par l'``intervalle'' ponctuel}
Le \deleted{\unsure{grou}} sous-\deleted{groupe}monoïde libre de $\operatorname{Fl}(\underline{\mathrm{Ty}})$ engendré par ces deux éléments \deleted{s'identifie \unsure{avec}}\inserted{est le} monoïde des flèches de dégénérescence \inserted{(pures)} quelconques,
